\clearpage
\fancyhead[L]{\small VOYAGE DE HUMBOLDT ET BONPLAND.}
\hypertarget{native-517-lower}{}\label{native-517-lower}
\begin{center}\small Göttingische gelehrte Anzeigen. Stück 91. S. 898..904. 10. Juni 1809.\\[3pt]\rule{.72\linewidth}{.3pt}\end{center}

\emph{Voyage} \textsc{de Humboldt} \emph{et} \textsc{Bonpland}. \emph{Quatrième Partie, Astronomie et Magnetisme. Premier Volume, contenant un recueil d'observations astronomiques, d'opérations trigonométriques, et de mesures barométriques, faites pendant le cours d'un voyage aux régions équinoxiales du Nouveau-Continent, depuis} 1799 \emph{jusqu'en} 1803. I. Lieferung 1808, II. Lieferung 1809. Großquart, 279 Seiten. Bei \textsc{Schöll} und \textsc{Cotta} in Paris und Tübingen.

Die Bedeutung und die erstaunliche Zahl der geografischen Ortsbestimmungen, die unser berühmter Landsmann \textsc{von Humboldt} auf seiner Reise durch die amerikanischen Tropenländer vorgenommen hat, würden für sich allein ausreichen, um dieser Reise einen der ersten Plätze unter den wissenschaftlichen Expeditionen in die Neue Welt zuzuweisen. Dass so ungemein viel \emph{von einem Mann} geleistet wurde und dass dies wiederum nur eines der Ziele war, die dieser außergewöhnliche Mann mit gleicher Begeisterung verfolgte, macht das Unternehmen einzigartig. Nichts war daher mehr zu wünschen, als dass der Schatz an Beobachtungen, den Herr \textsc{von Humboldt} nach Europa zurückgebracht hat, mit aller möglichen Gewissenhaftigkeit und Sorgfalt ausgewertet und genutzt würde, ohne flüch\SourceLeaf{520}tig und oberflächlich vorzugehen. Glücklicherweise fand Herr \textsc{von Humboldt} in Herrn \textsc{Oltmanns} einen Mann, der dieser umfangreichen und beschwerlichen Aufgabe gewachsen war und sie mit vorbildlicher Treue und ausdauerndem Fleiß auf sich nahm. Herr \textsc{Oltmanns} hatte diese rühmlichen Eigenschaften schon früher in anerkennenswerter Weise unter Beweis gestellt. Herr \textsc{von Humboldt} konnte diesem jungen Astronomen die Ableitung der Ergebnisse aus seinen Beobachtungen ohne Bedenken anvertrauen. Herr \textsc{von Humboldt} begnügte sich damit jedoch nicht. Damit alle die Güte seiner Beobachtungen, die Richtigkeit der Rechnungen von \textsc{Oltmanns} und die verhältnismäßige Zuverlässigkeit der Ergebnisse selbst beurteilen können, entschloss er sich, sämtliche Originalbeobachtungen ungekürzt neben den wichtigsten Schritten der Rechnungen von \textsc{Oltmanns} im Druck bekannt zu machen. Diese seltenen, für die Geografie Amerikas so wichtigen und unter so ungewöhnlichen Mühen und Gefahren errungenen Beobachtungen sind dieser Ehre gewiss sehr würdig. Allerdings sind \emph{die meisten} von ihnen nicht von der Art, dass man hoffen dürfte, ihnen künftig noch merklich genauere Ergebnisse zu entnehmen, als sich schon jetzt aus ihnen gewinnen lassen.

Über die Instrumente, die Herr \textsc{von Humboldt} auf seiner Reise mitführte, gibt er nach einer Mitteilung von Herrn \textsc{Oltmanns} in der Einleitung zum vorliegenden Werk selbst die nötige Auskunft. Wir wissen nicht, ob diese Einleitung bereits gedruckt ist. So entnehmen wir nur einigen beiläufigen Angaben, dass ein Chronometer von \textsc{Berthoud} und ein Sextant von \textsc{Ramsden} die am häufigsten gebrauchten Instrumente waren. Mitunter wurde ein kleiner zweizölliger Sextant von \textsc{Troughton} benutzt, und einige Male finden wir auch einen kleinen Quadranten von \textsc{Bird} erwähnt. Die Breitenbestimmungen beruhen am häufigsten auf Sonnenhöhen, seltener auf Sternhöhen. Die Längenbestimmungen wurden größtenteils mit dem Chronometer vorgenommen. An den wichtigeren Orten wurden aber auch Verfinsterungen der Jupitermonde und Monddistanzen beobachtet, in Cumana auch einmal eine Sonnenfinsternis. Die gemessenen Breiten können daher verständlicherweise nicht den Grad an Genauigkeit beanspruchen, den man mit fest aufgestellten Instrumenten oder mit Repetitionskreisen erreicht. Unter den Längenbestimmungen wird sich keine bis auf eine Bogenminute ganz verbürgen lassen. Das kann aber natürlich den Wert dieser Beobachtungen in einem Land nicht schmälern, in dem die Lage vieler der von Herrn \textsc{von Humboldt} bestimmten Orte bisher um mehrere Grad ungewiss war und die der meisten noch nie bestimmt worden war.

Es würde dem Zweck unserer Göttingischen gelehrten Anzeigen nicht entsprechen, unserem Reisenden Schritt für Schritt zu folgen und alle seine Bestimmungen ausführlich wiederzugeben. Deshalb begnügen wir uns damit, den Inhalt der drei Bücher, aus denen die beiden vorliegenden Lieferungen bestehen, kurz anzuzeigen. Das \emph{erste Buch} umfasst den Aufenthalt in Spanien und die Überfahrt nach Südamerika. Herr \textsc{von Humboldt} bestimmte dabei die Lage folgender Orte: Barcellona, Montserrat, Col de Balaquet, Venta de la Sienita, Valencia, die Ruinen von Sagunt, Madrid, Aranjuez, Corunna und St. Croix auf Teneriffa. Für einige dieser Orte liegen bereits gute Bestimmungen anderer Beobachter vor. Ihre ausgezeichnete Übereinstimmung mit den Bestimmungen des Herrn \textsc{von Humboldt} zeigt, wie viel Vertrauen diese verdienen. Herr \textsc{Oltmanns} hat hier und später an allen Orten, für die andere, ältere oder neuere Bestimmungen vorlagen, diese mit denen von \textsc{Humboldt} zusammengestellt, um sie anhand seiner Bestimmungen zu bestätigen oder zu berichtigen. Besonders wichtig und wertvoll sind diese Erörterungen bei den Orten, die als Hauptbezugspunkte aller in Südamerika vorgenommenen Längenbestimmungen gelten müssen, nämlich Cumana und Caraccas. Diese Untersuchungen machen deshalb einen beträchtlichen Teil des \emph{zweiten} und \emph{dritten Buches} aus. Aus einer sehr sorgfältigen Auswertung der Beobachtungen \textsc{Humboldt}s gewinnt \textsc{Oltmanns} für diese beiden wichtigen Orte folgende Ergebnisse:
\SourceLeaf{521}
\begin{center} \begin{tabular}{@{}lll@{}} \textit{Cumana}, Länge & $4^{\mathrm h}\,26^{\mathrm m}\,0^{\mathrm s}$ & Br. $10^\circ\,27'\,52''$ nördlich \\ \textit{Caraccas}, Länge & $4\quad 37\quad 40$ & Br. $10\quad 30\quad 50$ nördlich \\ \end{tabular} \end{center}
\noindent Den übrigen Teil des zweiten Buches füllen die Ortsbestimmungen für die Inseln Tabago und Trinidad, die die Mündungen des Orinoko und verschiedene Orte an der Küste von Neu-Andalusien. Der übrige Teil des dritten Buches behandelt die große Reise ins Innere von Südamerika. Auf ihr bestimmte Herr \textsc{von Humboldt} die Lage einer großen Zahl bemerkenswerter Orte am Orinoko und am Cassiquiare und stellte die Verbindung des ersteren über den letzteren und den Rio Negro mit dem Amazonas außer Zweifel. Dieser Teil der Reise \textsc{Humboldt}s ist umso wichtiger, als alle bisherigen Karten dieser Gegenden höchst fehlerhaft waren und zum Teil sogar die Breiten um mehr als 4 Grad falsch angaben.

Den Wert astronomischer Beobachtungen aus diesen Teilen Amerikas kann man erst richtig würdigen, wenn man die vielfältigen Schwierigkeiten bedenkt, mit denen der Beobachter dort zu kämpfen hat. In Caraccas ist das Wetter so unbeständig, dass es als höchst seltener Glücksfall gelten muss, wenn eine erwartete Beobachtung einmal gelingt. Herr \textsc{von Humboldt} hatte die Geduld, siebenundzwanzig Nächte durchzuwachen, um Verfinsterungen von Jupitermonden zu beobachten. Die meisten dieser Nächte waren ganz verloren. Zuweilen ist der Himmel 5 Minuten vor einer erwarteten Erscheinung noch völlig klar und im nächsten Augenblick bedeckt. Am 16. Dezember 1799, als sich ein Mond verfinstern sollte, hatte Herr \textsc{von Humboldt} den Planeten bei völlig klarem Himmel bereits im Gesichtsfeld des Fernrohrs. Sein Gehilfe \textsc{Bonpland} zählte schon die Sekunden. Nur noch eine Minute, und die Beobachtung wäre gemacht gewesen. Doch plötzlich entsteht ein so dichter Nebel, dass selbst vom Planeten nichts mehr zu sehen ist. In der Wildnis des Orinoko ist die Lage des Beobachters noch außergewöhnlicher. Während der Nacht, schreibt Herr \textsc{von Humboldt}, erschreckte uns ringsum das Gebrüll der Tiger, begleitet vom Geheul unserer Hunde. Die Krokodile, angelockt von dem Feuer, das unsere indigenen Führer am Ufer angezündet hatten, richteten sich zur Hälfte aus dem Wasser auf, als wollten sie uns beobachten. Aber selbst solche Szenen lassen einen mit der Zeit gleichgültig, wenn man sich erst daran gewöhnt hat.

Dem zweiten Buch ist unter dem Titel \emph{Essai sur les refractions dans la zone torride} noch eine Abhandlung von Herrn \textsc{von Humboldt} angehängt. Darin betrachtet er aus verschiedenen, zum Teil neuen Blickwinkeln die interessante Frage, ob in unterschiedlichen Klimaten dieselbe Refraktionstafel angewandt werden darf. Herr \textsc{von Humboldt} betrachtet zuerst den Einfluss der unterschiedlichen Mischungsverhältnisse der Bestandteile der atmosphärischen Luft und ihres Feuchtigkeitszustands. Er zeigt, dass kein Grund besteht, daraus eine unterschiedliche Refraktion für verschiedene Klimate zu vermuten. Wichtiger ist, wie das Gesetz der Temperaturabnahme in den höheren Luftschichten die Refraktion bei kleinen Höhen und am Horizont beeinflusst. Die Frage, ob dieses Gesetz in den verschiedenen Erdzonen dasselbe ist, ist daher von größter Wichtigkeit. Herr \textsc{von Humboldt} hat dazu mehrere sehr bemerkenswerte direkte Versuche angestellt, die hinsichtlich der Temperaturabnahme in der heißen Zone überraschend gut übereinstimmen. Aus allen ergibt sich ein Höhenunterschied von 122 Toisen für eine Temperaturabnahme um 1 Grad Réaumur. Das unterscheidet sich fast gar nicht von dem, was \textsc{Gay-Lussac} in einem Ballon in einer Höhe von 6980 Metern über Paris feststellte. Es scheint also, dass dieses Gesetz in der warmen Jahreszeit in der gemäßigten und der heißen Zone genau dasselbe ist. Bei niedrigerem Thermometerstand dagegen scheint die Temperatur mit der Höhe langsamer abzunehmen. In Ermangelung direkter Beobachtungen hierzu fügt Herr \textsc{von Humboldt} die Ergebnisse bei, die Herr \textsc{Mathieu} aus zwei bekannten, in Torneå bei sehr starker Kälte ange\SourceLeaf{522}stellten Beobachtungen mithilfe der Hypothese von \textsc{Laplace} über die Abnahme der Luftdichte abgeleitet hat. Für eine Temperaturabnahme um 1 Grad Réaumur ergeben sie einen Höhenunterschied von $156\frac{1}{2}$ Toisen. Nach diesem Ergebnis sollte man also meinen, dass die Refraktionen in der kalten Zone eine andere Tafel als in der gemäßigten erfordern, die in der heißen Zone dagegen nicht. Das ist gerade das Gegenteil von dem, was einerseits \textsc{Bouguer} und andererseits \textsc{Maupertuis} und \textsc{Lemonnier} behauptet haben. Man weiß jedoch, dass \textsc{Bouguer} seine Beobachtungen selbst geändert hat, um sie mit den vermeintlich kleineren Refraktionen in Einklang zu bringen. Dagegen sprechen ferner die Beobachtungen von \textsc{Legentil}, die ebenfalls in der heißen Zone, in Pondichery, angestellt und von \textsc{Delambre} neu berechnet wurden. Sie vertragen sich recht gut mit der Tafel von \textsc{Bradley}, aber nicht mit der von \textsc{Bouguer}. Schließlich hat Herr \textsc{von Humboldt} selbst mehrere Beobachtungen zur Refraktion in Cumana, Caraccas, Acapulco und auf der Insel Cuba angestellt, die durchweg größere Zahlen ergeben als die Tafel von \textsc{Bouguer}. Aus all diesen Gründen scheint es am besten, die in der gemäßigten Zone aufgestellte Refraktionstafel auch in der heißen Zone beizubehalten, bis ein geübter Beobachter dort mit Geräten, die der neueren Beobachtungskunst entsprechen, eine Reihe zahlreicher Beobachtungen unter den verschiedensten atmosphärischen Bedingungen angestellt haben wird. In der gemäßigten und kalten Zone dagegen nimmt die Temperatur in den höheren Luftschichten umso langsamer ab, je niedriger der Thermometerstand in der untersten Luftschicht ist. Das wird es nötig machen, die Art zu ändern, wie unsere Refraktionstafeln die Temperatur berücksichtigen. Leider! Die Hoffnung, hier je zu einer Klärung zu kommen, wird durch die Bemerkung \textsc{Delambre}s sehr gedämpft, dass sich die Refraktion am Horizont bei fast unveränderten Barometer- und Thermometerständen häufig um 4 Minuten ändert, ohne dass man einen Grund dafür anzugeben weiß.

Wir müssen noch bemerken, dass der Druck dieses Werks von \textsc{Humboldt} eine äußere Eleganz aufweist, die seinem inneren Wert entspricht. Wir wünschten, auch die Genauigkeit des Drucks selbs bei den Zahlenangaben so loben zu können. Doch haben wir ungern bemerkt, dass die Korrektur in dieser Hinsicht ziemlich nachlässig ausgeführt wurde. Da der Abdruck von Originalbeobachtungen nur so weit von Wert ist, wie man allen Zahlen vertrauen kann, ist zu erwarten, dass der Herausgeber sie noch einmal sorgfältig mit der Handschrift vergleicht und alle Druckfehler am Ende gewissenhaft angibt.
